
\documentclass[../../ps_q.tex]{subfiles}

\begin{document}

\section{静電場}

\subsection{点電荷が作る電場}
    \subfile{01_point-charge/ps_em_ef_01_q.tex}

    \vspace{2cm}
    \vfill

\subsection{平行一様電場・電位の基準}
    \subfile{02_uniform/ps_em_ef_02_q.tex}

    \vspace{2cm}
    \vfill

\subsection{拡がった電荷分布の作る電場}
    \subfile{03_continuous-charge/ps_em_ef_03_q.tex}

    \vspace{2cm}
    \vfill

\subsection{電位から電場を逆算する①}
    \subfile{04_from-potential-1/ps_em_ef_04_q.tex}

    \vspace{2cm}
    \vfill

\subsection{電位から電場を逆算する②}
    \subfile{05_from-potential-2/ps_em_ef_05_q.tex}
    
    \vspace{2cm}
    \vfill

\subsection{※点電荷が作る電場}
    \subfile{06_add-point-charge/ps_em_ef_06_q.tex}

    \vspace{2cm}
    \vfill

\subsection{※平行一様電場}
    \subfile{07_add-uniform/ps_em_ef_07_q.tex}

    \vspace{2cm}
    \vfill

\subsection{※電気力線}
    \subfile{08_add-field-lines/ps_em_ef_08_q.tex}

    \vspace{2cm}
    \vfill

\subsection{※等電位線}
    \subfile{09_add-equipotential-lines/ps_em_ef_09_q.tex}

    \vspace{2cm}
    \vfill

\end{document}
